\title{Parameter-Efficient Orthogonal Finetuning via Factorization}

\begin{abstract}
The prevalence of large foundation models has surged, but the immense computational resources needed to train such models from the ground up often make this impractical. Consequently, methods for effectively tailoring these robust models to specific downstream applications are gaining significance. In this work, we investigate Orthogonal Finetuning (OFT)—a theoretically grounded finetuning scheme—for adapting pre-trained models to new tasks. While OFT is known for its strong generalization properties, it nevertheless involves tuning a relatively large parameter set, owing to the inherent high dimensionality of orthogonal matrices. To tackle this limitation, we analyze OFT through the perspective of information transfer, identifying essential criteria for improved parameter-efficiency. Drawing inspiration from the efficiency of the Cooley-Tukey fast Fourier transform in facilitating rapid information propagation, we introduce a compact orthogonal parameterization utilizing butterfly network structures. We embed this parameterization within OFT to formulate Orthogonal Butterfly~(BOFT), a new finetuning technique optimized for parameter efficiency. BOFT generalizes OFT by treating it as a particular instance within a broader orthogonal finetuning framework. Extensive empirical evaluation demonstrates the effectiveness of this method for adapting a range of models—including large vision transformers, language models, and text-to-image diffusion architectures—to diverse tasks across vision and language domains.
\end{abstract}

\section{Introduction}

Groundbreaking systems such as ChatGPT~\cite{brown2020language,chen2021evaluating} and Stable Diffusion~\cite{rombach2022high} underscore the strong generalization capability intrinsic to large foundation models. The dramatic escalation in model capability has gone hand in hand with a surge in parameter counts

(\emph{e.g.}, GPT-3 boasts approximately 175B parameters). As model sizes increase, training foundation models ab initio is becoming increasingly prohibitive for most research teams. Therefore, advancing the field hinges on methodologies that enable adapting pre-trained foundation models without restarting the training process. This necessitates mechanisms for efficiently specializing these models for new downstream objectives. Three predominant strategies exist for efficient adaptation: \emph{model finetuning}~\cite{hulora2022,zaken2022bitfit,guo2020parameter,raffel2020exploring,qiu2023controlling,zhang2022adaptive,chavan2023one}, where the underlying model remains intact and only select parameters are adjusted; \emph{adapter tuning}~\cite{rebuffi2017learning,houlsby2019parameter,pfeiffer2020adapterfusion,lin2020exploring,he2021towards}, which inserts extra trainable components into the original architecture; and \emph{prompt tuning}~\cite{li2021prefix,lester2021power}, whereby auxiliary learnable prefixes are prepended to the model inputs. Among these, model finetuning is distinguished by its simplicity and efficacy, crucially achieving adaptation with zero additional inference delay.

The core idea underlying model finetuning is to maintain a close resemblance between the pretrained and finetuned models, as measured by certain criteria, thereby retaining the knowledge acquired during pretraining. Conventional finetuning techniques frequently utilize a reduced learning rate; this practical measure implicitly limits the deviation between the pretrained and finetuned models. Recognizing the organized structure inherent to weight matrices, more theoretically motivated strategies focus on preserving the relational properties of these matrices—specifically, the pairwise angular relationships among neurons. This perspective has given rise to Orthogonal Finetuning~(OFT)~\cite{qiu2023controlling}, a new paradigm in model finetuning, which operates by applying a uniform orthogonal transformation to all neurons within the same layer. This approach mathematically guarantees that the pairwise angles between neurons are conserved throughout the finetuning stage. While OFT has achieved notable results in terms of both generalization and optimization when employed on text-to-image diffusion models~\cite{qiu2023controlling}, a practical challenge emerges: the orthogonal matrices involved are high-dimensional, leading to a considerable increase in trainable parameters. In response, OFT employs a block-diagonal architecture to mitigate the parameter count. Nonetheless, this gain in parameter efficiency introduces trade-offs—the block-diagonal orthogonal matrix imposes a rigid sparsity structure and applies transformations independently across its constituent blocks. Such a constraint may lead to unwelcome inductive biases (for example, limiting expressiveness and preventing the orthogonal matrix from representing certain canonical linear transformations), and, crucially, the problem of discovering an optimal sparsity pattern for these blocks remains unsolved.

The central challenge here lies in constructing a dense orthogonal matrix without incurring excessive parameter costs. Although it initially appears impractical—since a dense, $d$-dimensional orthogonal matrix normally needs $\mathcal{O}(d^2)$ parameters—we introduce a new strategy: building the dense orthogonal matrix by sequentially composing several sparse, factorized matrices. This methodology is informed by the realization that we can achieve parameter reduction by trading additional computation per iteration for decreased storage. In this scheme, the orthogonal matrix is realized as a product of sparse matrices, necessitating their repeated multiplication throughout training. To contextualize this factorization, we reframe dense orthogonal matrix generation as an information transmission challenge. In this view, synthesizing a dense orthogonal matrix via sparse matrix products becomes analogous to disseminating information across a grid-like network. This information transmission perspective allows for a range of sparse matrix factorizations that are sufficiently expressive to produce dense matrices, yet require substantially fewer trainable parameters.

To reach parameter efficiency within this framework, we are inspired by the butterfly network structure underlying the Cooley-Tukey fast Fourier transform algorithm~\cite{cooley1965algorithm}, which enables rapid data communication between nodes~\cite{klasing1994broadcasting}. In particular, the butterfly topology intrinsic to the Cooley-Tukey method suggests an approach known as butterfly factorization for decomposing matrices sparsely. Through butterfly factorization, a dense $d$-dimensional matrix can be expressed as a product of $\mathcal{O}(\log d)$ sparse matrices, each containing only $\mathcal{O}(d)$ nonzero elements. Enforcing orthogonality at the level of each sparse factor leads to an effective orthogonal parameterization utilizing merely $\mathcal{O}(d\log d)$ parameters—a dramatic reduction from conventional parameter counts. Building upon this compact orthogonal representation, we introduce Orthogonal Butterfly~(BOFT), a parameter-efficient finetuning approach. BOFT generalizes OFT as a particular instance and serves as a broad framework for orthogonal finetuning. Notably, both block-diagonal and butterfly configurations exploit sparsity to reduce trainable parameters within orthogonal matrices. A comparison with LoRA~\cite{hulora2022} is informative: both low-rank and sparse matrices are principal classes of structured matrices~\cite{candes2011robust} that enable greater parameter efficiency.

In contrast to the block-diagonal design leveraged by OFT to balance model expressiveness and regularity, BOFT employs the butterfly structure to achieve a more seamless transition between matrices that span the entire orthogonal group (allowing high expressivity) and those close to the identity (encouraging regularity). This approach allows us to identify a more compact subset of hypotheses within the orthogonal group that are well-suited for specific downstream tasks. Notably, the butterfly structure has been widely adopted in efficient implementations of various linear transforms, including the discrete Fourier and discrete cosine transforms. We suggest that this structurally constrained parameterization may act as an inherent inductive bias, fostering better generalization and offering resistance against overfitting. Our reasoning is grounded in the observation that butterfly structures can readily reproduce numerous classical linear transforms, a capability absent in the block-diagonal framework of OFT. We summarize our contributions as follows:

\begin{itemize}[leftmargin=*,nosep]
    \item We address parameter efficiency in orthogonal finetuning through a novel perspective rooted in information transmission. Specifically, we reinterpret the search for compact, dense orthogonal matrices as an information flow problem mapped onto a grid-structured graph.
    \item Drawing inspiration from the butterfly architectures foundational to the Cooley-Tukey algorithm, we introduce Orthogonal Butterfly—an orthogonal finetuning approach that is both parameter-efficient and situated within our information transmission framework.
    \item We supply several theoretical justifications explaining how BOFT can achieve a drastic reduction in trainable parameters while retaining robust expressivity and stable optimization. The matrix factorization in BOFT also naturally enables an appealing form of weight interpolation (refer to Figure~\ref{fig:boft_interpolation}).
    \item For the first time, we extend orthogonal finetuning to a broad range of tasks beyond controllable text-to-image synthesis~\cite{qiu2023controlling}, thereby demonstrating its viability as a versatile approach for model adaptation. We deploy BOFT across diverse application domains—including both computer vision and natural language processing—and show that it significantly surpasses existing leading methods, confirming its parameter efficiency and superior generalization.
\end{itemize}

\section{Related Work}

\textbf{Parameter-efficient finetuning (PEFT)}. As the scale and capabilities of foundation models (\emph{e.g.}, \cite{brown2020language,radford2021learning,rombach2022high,kirillov2023segment}) continue to increase, the challenge of finetuning these large models for downstream applications in a way that uses parameters efficiently has received growing attention~\cite{treviso2023efficient,ding2023parameter,lialin2023scaling}. Of the numerous techniques introduced under the umbrella of PEFT~\cite{houlsby2019parameter,aghajanyan2020intrinsic,hulora2022,edalati2022krona,wang2022adamix,gheini2021cross,zaken2022bitfit,guo2020parameter,sung2021training,ansell2022composable,lester2021power,li2021prefix,vu2022spot,he2021towards,mao2021unipelt,karimi2021compacter,liu2022few,sung2022lst,chen2022parameter,jia2022visual,chen2022adaptformer,zhang2022neural,jie2023fact,lian2022scaling,luo2023towards}, the family of reparameterization-based approaches~\cite{aghajanyan2020intrinsic,hulora2022,edalati2022krona} bears the closest relation to our present study. LoRA~\cite{hulora2022}, for instance, adapts pretrained weight matrices by appending the product of two trainable low-rank matrices, and has demonstrated strong effectiveness on natural language processing tasks. Recognizing that LoRA constrains the rank of inserted matrices to a fixed value across all layers, recent works~\cite{zhang2022adaptive,valipour2022dylora,zhang2023increlora} introduce approaches that flexibly vary rank per layer to better distribute the limited parameter resources. The widespread adoption of such low-rank parameterizations is largely attributable to their simplicity and efficiency~\cite{dettmers2023qlora,zi2023delta,chavan2023one}. Drawing inspiration from theoretical insights regarding the role of hyperspherical energy in generalization~\cite{liu2018learning,liu2021orthogonal}, \cite{qiu2023controlling} developed orthogonal finetuning (OFT), presenting an alternative approach in which each layer’s neurons are linearly transformed by a learned orthogonal matrix; this method offers both improved generalization and more stable training than LoRA in the context of text-to-image diffusion models. Nonetheless, a drawback of OFT is its typically larger number of trainable parameters relative to LoRA, motivating the need to increase its parameter efficiency. Additionally, the potential of OFT across a broader range of adaptation problems (apart from text-to-image diffusion, as in~\cite{qiu2023controlling}) has not been established. The BOFT method leverages butterfly factorization to render OFT more parameter-efficient, which enables us to extend and empirically verify the advantages of orthogonal finetuning in diverse adaptation settings.

\textbf{Butterfly structures}. The radix-2 Cooley-Tukey algorithm offers an efficient way to compute the $N$-point discrete Fourier transform by recursively breaking it down into two separate $\frac{N}{2}$-point DFTs. This procedure introduces a so-called butterfly structure, formally represented as the product of a sequence of sparse matrices, collectively termed a butterfly matrix. The butterfly matrix has previously been employed to parametrize orthogonal matrices as a means to forgo pivoting in Gaussian elimination and to improve computational efficiency~\cite{parker1995random}, to facilitate stable training in recurrent neural networks~\cite{jing2017tunable}, and for kernel approximation~\cite{munkhoeva2018quadrature}. Several works~\cite{dao2019learning,dao2019kaleidoscope} have explored learning rapid linear transformations through butterfly-based parametrizations, while others~\cite{chen2022pixelated,dao2022monarch} exploit these matrices for the efficient training of deep networks. Beyond these applications, butterfly structures have proved effective in tasks like fast matrix-vector computations~\cite{michielssen1996multilevel,o2010algorithm}, compact matrix representations~\cite{li2015butterfly}, and information transfer in networks~\cite{klasing1994broadcasting,li2003linear,rajasingh2016transmission}. In this study, unlike previous usage, we concentrate on utilizing butterfly structures specifically to augment the parameter efficiency of OFT.

\section{An Information Transmission View on Orthogonal Finetuning}

We begin by outlining foundational aspects of OFT. For adapting a pretrained weight matrix, OFT employs a reparameterization wherein the updated weight matrix is represented as the product of a trainable orthogonal matrix and the fixed pretrained weights. In contrast to LoRA, which introduces changes through the addition of a low-rank matrix, OFT relies on multiplicative updates via an orthogonal transformation. LoRA achieves parameter efficiency by imposing a low-rank constraint, whereas the standard OFT framework leverages a block-diagonal structure on the orthogonal matrix~\cite{qiu2023controlling}; parameter efficiency increases as the block size decreases. This distinction is illustrated in Figure~\ref{fig:comp_lora}. The principal motivation behind utilizing an orthogonal transformation for fine-tuning lies in its capacity to conserve pairwise neuronal angles~\cite{liu2018learning,liu2021orthogonal,qiu2023controlling}, thereby retaining much of the semantic information acquired during pretraining.

Formally, OFT optimizes an orthogonal matrix $\bm{R}\in\mathbb{R}^{d\times d}$ associated with a pretrained linear weight matrix $\bm{W}^0\in\mathbb{R}^{d\times n}$. This changes the computation in the forward pass from $\bm{z} = (\bm{W}^0)^\top \bm{x}$ to $\bm{z} = (\bm{R}\bm{W}^0)^\top \bm{x}$, where $\bm{x}\in\mathbb{R}^{d}$ and $\bm{z}\in\mathbb{R}^{n}$ denote the input and output vectors, respectively. Zero initialization is achieved by setting $\bm{R}$ initially as the identity matrix. To guarantee that $\bm{R}$ remains orthogonal throughout training, we adopt Cayley parameterization following~\cite{qiu2023controlling,liu2021orthogonal}: $\bm{R} = (\bm{I} + \bm{Q})(\bm{I} - \bm{Q})^{-1}$, where $\bm{Q}$ is a skew-symmetric matrix ($\bm{Q} = -\bm{Q}^\top$). For the purpose of parameter reduction, OFT expresses $\bm{R}$ as a block-diagonal matrix, $\text{diag}(\bm{R}_1, \bm{R}_2, \ldots, \bm{R}_r)$, with each block $\bm{R}_i \in \mathcal{R}^{b\times b}$, and $br = d$. 

This block-diagonal design ensures fewer parameters but introduces an artificial grouping assumption: the dimensions of each neuron's vector (columns of $\bm{W}^0$) are partitioned into $r$ groups, and each group is independently transformed via a separate orthogonal block. Despite empirical validation of this method, such arbitrary dimension partitioning lacks theoretical justification, thus motivating interest in dense orthogonal matrices. Consequently, the following question naturally arises: \emph{Is it feasible to construct a dense orthogonal matrix while retaining the same parameter-efficiency?}

To tackle this issue, we suggest constructing a dense orthogonal matrix by multiplying several sparse orthogonal matrices. To systematically analyze the sparsity configurations that arise in orthogonal matrix decompositions, we reinterpret the task of forming a dense orthogonal matrix within the OFT framework as an information transmission process. Concretely, forming a dense matrix $\bm{R}\in\mathbb{R}^{d\times d}$ through a sequence of $m$ square matrices, such that $\thickmuskip=2mu \medmuskip=2mu \bm{R}=\bm{B}_m\bm{B}_{m-1}\cdots\bm{B}_1$, is analogous to propagating information through a lattice comprising $d\times (m+1)$ nodes, as depicted in Figure~\ref{fig:info_trans}. This information transmission analogy arises from perceiving a dense $d$-by-$d$ matrix as a fully-connected mapping from $d$ input nodes to $d$ output nodes. In this setup, a zero entry $R_{ij}$ in $\bm{R}$ signifies that the $j$-th input node cannot transmit information to the $i$-th output node, while a nonzero $R_{ij}$ indicates possible information flow. As such, expressing $\bm{R}$ as a sequence of matrix products, $\bm{B}_m\bm{B}_{m-1}\cdots\bm{B}_1$, can be equivalently viewed as enacting iterative exchanges of information, dictated by the graph structures underlying each $\bm{B}_i$. Here, information flows first via $\bm{B}_1$, culminating in $\bm{B}_n$. For a concrete illustration, Figure~\ref{fig:info_trans} demonstrates the decomposition $\bm{R}=\bm{B}_5\bm{B}_4\bm{B}_3\bm{B}_2\bm{B}_1$, showing both the associated sparsity structures and the corresponding graph. This graphical representation results from unfolding the matrix multiplication process. Within this framework, $\bm{B}_i$ acts as a connectivity matrix, linking the $i$-th layer of nodes to those at level $(i+1)$. More precisely, entry $(j_1,j_2)$ in $\bm{B}_i$ encodes whether a directed edge exists from node $j_2$ in layer $i$ to node $j_1$ in layer $i+1$, where a zero value denotes the absence of such an edge. Ensuring $\bm{B}_5\bm{B}_4\bm{B}_3\bm{B}_2\bm{B}_1$ is dense thus requires every node at the initial level to transmit information to each node at the sixth level. If we examine $\bm{R}=\bm{B}_4\bm{B}_3\bm{B}_2\bm{B}_1$, connecting source nodes in the first level to receivers at the fifth, it becomes apparent that node 1 cannot reach node 3, implying that the $(3,1)$ element is zero in the associated sparsity pattern. Similar correspondences between induced graph connectivity and sparsity arise in the cases of $\bm{R}=\bm{B}_3\bm{B}_2\bm{B}_1$ and $\bm{R}=\bm{B}_2\bm{B}_1$. Broadly, for any dense matrix $\bm{R}\in\mathbb{R}^{d\times d}$, the constituent factors $\bm{B}_m,\ldots,\bm{B}_1$ must collectively specify a set of directed edges on a $d\times(m+1)$ lattice, with each edge linking nodes between consecutive layers (i.e., adjacent columns), in such a manner that each source node at level one can transmit information to all nodes at level $(m+1)$.

Viewing OFT through the lens of information transmission yields insight into the sparse patterns seen in its factorization matrices. Figure~\ref{fig:blk_diag} illustrates how the original OFT yields a block-diagonal form for $\bm{R}$. This structure, while effective at parameter reduction, inherently restricts $\bm{R}$ from being dense. The main objective becomes how to assemble a dense orthogonal matrix by composing $m$ sparse orthogonal matrices, aiming to minimize the number of meaningful, learnable parameters.

From the information transmission standpoint, two primary criteria emerge: \emph{(i)} \textbf{dense connectivity}, where each node at the input layer is connected to every output node by at least one path, and \emph{(ii)} \textbf{minimal free edges}, i.e., the total edge count is minimized while preserving orthogonality. The orthogonality condition introduces specific constraints on inter-level connections. To ensure that each matrix $\bm{B}_i$ retains full rank—a prerequisite for orthogonality—$d$ edges are necessary to create a one-to-one mapping between nodes in adjacent layers. Thus, at least $d$ edges must be present between neighboring levels (for example, in Figure~\ref{fig:info_trans}, there are four such edges). Importantly, these $d$ connections are mandated by orthogonality and should be excluded from the count of trainable edges, since they represent fixed (non-trainable) entries—typically restricted to $\pm 1$ for $d\times d$ orthogonal matrices with only $d$ non-zero entries.

Given that orthogonal matrices have fewer independent parameters than dense matrices, this constraint introduces interdependencies among edges. For instance, in the case of a $2\times 2$ orthogonal matrix, if the $(1,1)$ entry is zero, orthogonality dictates that the $(2,2)$ entry must also be zero—demonstrating that omitting one edge could necessitate omitting another. Hence, orthogonality can cause specific patterns of non-zero entries (edges) to either co-occur or be jointly excluded. Consequently, the information transmission framework becomes especially useful in analyzing OFT: it focuses attention on the trainable non-zero entries of $\bm{R}$ and their configuration, independent of their precise numerical values.

A straightforward approach to establishing dense connections between two levels requires $\mathcal{O}(d^2)$ edges, corresponding to a single dense orthogonal matrix, which actually involves $d^2 - d$ edges (for instance, when $d=4$, there are 12 edges). Figure~\ref{fig:info_trans} demonstrates an example utilizing matrix factorization that uses just $10$ edges in total, fewer than what is required by a single dense orthogonal matrix. This perspective allows for an analysis of OFT’s parameter efficiency from a graphical viewpoint, making it straightforward to devise viable factorizations within this setting. Our methodology is partly inspired by the so-called butterfly graph topology from the Cooley-Tukey algorithm~\cite{cooley1965algorithm}, a structure that achieves full connectivity between $d$ input and $d$ output nodes with a mere $\mathcal{O}(d\log d)$ edges. As illustrated in Figure~\ref{fig:info_trans}, the employed topology achieves dense connectivity using $10$ edges, whereas a butterfly network achieves the same with only $8$ edges. In what follows, we explain how leveraging the butterfly architecture can further enhance parameter efficiency.

\section{Orthogonal Parameterization by Butterfly Factorization}

Originally introduced in the context of the Cooley-Tukey algorithm to accelerate the fast Fourier transform, the butterfly structure has the remarkable property that a localized perturbation in the frequency domain translates to a global effect in the spatial domain—a principle closely mirroring our challenge of information propagation, where all nodes at the input level can relay data to every node at the output. Furthermore, the butterfly structure has found widespread application as a network topology for efficient data transmission in computing systems~\cite{solihin2015fundamentals,li2011linear}. Assuming $k\geq 2$ and $k$ is a power of two, we define the \emph{butterfly factor} $\bm{B}^F(k)$ by

\begin{equation}\label{eq:but_factor}
\begin{aligned}
    \bm{B}^F(k)=\begin{bmatrix}
    \text{diag}(\bm{d}_1) & \text{diag}(\bm{d}_2) \\
    \text{diag}(\bm{d}_3) & \text{diag}(\bm{d}_4)
    \end{bmatrix}\in\mathbb{R}^{k\times k},
\end{aligned}
\end{equation}

where each $\bm{d}_i\in\mathbb{R}^{\frac{k}{2}},\forall i$ is a vector. For $d=2^N$, the $d$-dimensional \emph{butterfly matrix} $\bm{B}(d)\in\mathbb{R}^{d\times d}$ is then recursively defined as

\begin{equation}
\begin{aligned}
    \!\!\bm{B}(d)=\tilde{\bm{B}}(d,d)\!\cdot\!\begin{bmatrix}
    \bm{B}_1(\frac{d}{2})\!\!\!\!\! & \bm{0} \\
    \bm{0}\!\!\!\!\! &\bm{B}_2(\frac{d}{2})
    \end{bmatrix}= \tilde{\bm{B}}(d,d)\tilde{\bm{B}}(d,\frac{d}{2})\cdots\tilde{\bm{B}}(d,2),
\end{aligned}
\end{equation}

Here, $\bm{B}_1(\frac{d}{2})$ and $\bm{B}_2(\frac{d}{2})$ represent two butterfly matrices, each with dimensions $\frac{d}{2}$. The \emph{butterfly component} is then constructed as $\thickmuskip=2mu \medmuskip=2mu \tilde{\bm{B}}(d,k)=\text{diag}(\bm{B}^F_1(k),\cdots,\bm{B}^F_{d/k}(k))$, forming a block-diagonal matrix of size $d \times d$, where each block has dimension $k$. The blocks $\bm{B}^F_i(k)$ for all $i$ correspond to the butterfly factors specified in Equation~\ref{eq:but_factor}. With this setup, the butterfly matrix can be employed to parameterize an orthogonal matrix by ensuring all constituent matrices $\tilde{\bm{B}}(d,k)$ in $\bm{B}(d)$ are orthogonal. Focusing initially on the case of the block-diagonal matrix $\tilde{\bm{B}}(d,2)$ with block size 2, we can ensure its orthogonality by parameterizing every $2 \times 2$ block using a Cayley transform or $2$-dimensional rotation, following the method in \cite{liu2021orthogonal,qiu2023controlling}. Since the sparsity structure of any butterfly component is a permutation of the pattern in $\tilde{\bm{B}}(d,2)$, the same strategy facilitates orthogonal parameterization for all butterfly components, resulting in an efficient representation of orthogonal matrices using multiple $2\times 2$ orthogonal blocks. Extending the construction in \cite{chen2022pixelated}, we define the block butterfly component $\tilde{\bm{B}}^b(d,k)$ so that each entry in $\bm{d}_i$ for every $i$ is now a $b\times b$ matrix. To maintain orthogonality in the block butterfly component $\tilde{\bm{B}}^b(d,2)$, each $2b\times 2b$ block is enforced to be orthogonal. For butterfly components with $k>2$, i.e., $\tilde{\bm{B}}^b(d,k)$, their sparsity structures can be interpreted as block-wise permutations of that in $\tilde{\bm{B}}^b(d,2)$ and thus are parameterized as orthogonal matrices analogously. Synthesizing these components yields the BOFT forward computation as follows:

\begin{equation}
    \begin{aligned}\nonumber
    \bm{z} = \big(\bm{R}(m,b)\cdot\bm{W}^0\big)^\top \bm{x},~~~~\text{subject to}~\bigg\{\bm{R}(m,b) = \prod_{i=1}^m \tilde{\bm{B}}^b_{(d,i)}~~\&~~\underbrace{\big(\tilde{\bm{B}}^b_{(d,j)}\big)^\top\tilde{\bm{B}}^b_{(d,j)} = \tilde{\bm{B}}^b_{(d,j)}\big(\tilde{\bm{B}}^b_{(d,j)}\big)^\top = \bm{I}_d}_{\forall j \in [1, m]}\bigg\},
    \end{aligned}
\end{equation}

For clarity, we use $\tilde{\bm{B}}^b(d,2^{m-i+1})$'s shorthand $\tilde{\bm{B}}^b_{(d,i)}$, and $\bm{I}_d$ stands for the $d$-dimensional identity matrix. The matrix $\bm{R}(m,b)\in\mathbb{R}^{d\times d}$ is an orthogonal transformation formed as the product of multiple orthogonal butterfly matrices. For ease of reference, we write BOFT using $\bm{R}(m,\frac{b}{2})$ as BOFT($m$,$b$), taking $b \geq 2$. When setting $m=1$, BOFT($1$,$b$) reduces to the block-diagonal OFT~\cite{qiu2023controlling} where each block is of size $b$. At the extreme, BOFT($1$,$d$) is equivalent to the classical OFT~\cite{qiu2023controlling} involving a dense orthogonal matrix without block constraints. Choosing $m=\log \frac{2d}{b}$ and $b$ accordingly leads BOFT to utilize a block butterfly matrix $\bm{B}^{\frac{b}{2}}(d)$ for $\bm{R}$, resulting in a dense orthogonal matrix construction. More broadly, BOFT($m$,$b$) requires a total of $\frac{1}{2}(b-1)dm$ learnable parameters for updating a linear layer of dimension $d\times n$. If the butterfly matrix configuration is chosen (\emph{i.e.}, $m = \log d$, $b = 2$), then BOFT only requires $\mathcal{O}(d\log d)$ parameters. In contrast, an unconstrained OFT using a full dense orthogonal matrix needs $\mathcal{O}(d^2)$ parameters, and a block-diagonal OFT with $r$ blocks needs $\mathcal{O}(bd)$. As such, generating a dense orthogonal matrix via OFT requires $b=d$, but BOFT achieves this flexibility for any value of $b$.

\textbf{BOFT Identity Initialization}. Most finetuning approaches begin from the precise state of a pretrained model to avoid significant deviation from its learned knowledge. As a case in point, LoRA applies zero initialization to its low-rank weight parameters. Analogously, BOFT initializes each butterfly matrix with the identity transformation (specifically, skew-symmetric components are initialized with zeros when using the Cayley transform).

\textbf{Multiplicative Dropout in BOFT}. LoRA~\cite{hulora2022} incorporates Dropout layers with the low-rank adaptation to control overfitting risk. While standard Dropout~\cite{srivastava2014dropout} is well-suited for LoRA, it is ill-posed for BOFT due to BOFT's use of multiplicative updates. We thus introduce a multiplicative Dropout approach tailored for BOFT. Since $\bm{R}(m,b)$ is constructed from $m$ orthogonal butterfly matrices, each capable of being permuted into $2b\times 2b$ block-diagonal form, our proposed Dropout operates by randomly selecting a proportion $p_1$ of butterfly matrices and $p_2$ fraction of the diagonal blocks per butterfly. These randomly chosen blocks are then replaced with identity matrices during training.

\section{Intriguing Insights and Discussions}

\textbf{Expressivity of BOFT}. While the butterfly structure combined with permutations can exactly represent many classical fast linear transforms~\cite{dao2019learning,dao2019kaleidoscope} (\emph{e.g.}, fast Fourier transform, Hadamard transform), it remains unclear to what degree our orthogonal butterfly matrix can approximate arbitrary orthogonal matrices. To examine this, we perform a simulation task where a randomly generated dense orthogonal matrix~\cite{anderson1987generation} of size $512\times 512$ is approximated. As shown in Figure~\ref{fig:para_eff}, outcomes are averaged over 10 distinct random seeds. Here, the y-axis measures the approximation error, while the x-axis corresponds to the number of trainable parameters in use. Each curve of identical color represents a particular BOFT block size, and the leftmost marker in each curve denotes the error for BOFT($1$,$b$) (\emph{i.e.}, a vanilla block-diagonal OFT with block size $b$). Across configurations, BOFT typically offers improved parameter efficiency over OFT. For instance, BOFT($9$,$2$) provides greater expressive capability than BOFT($1$,$16$), despite utilizing significantly fewer parameters. Using a smaller $b$ with a larger $m$ often results in superior parameter efficiency for BOFT; to illustrate, BOFT($6$,$4$) achieves a comparable approximation error to BOFT($2$,$16$) with substantially fewer parameters. Overall, the butterfly matrix forms a more structured subset within the orthogonal group than a block-diagonal design, enabling BOFT to be provably more expressive than OFT given the same block size.

\begin{theorem}[Expressivity of BOFT]\label{thm:exp_boft}
    BOFT surpasses OFT in expressiveness for an equivalent block size. Any orthogonal matrix of size $d$ can be constructed by multiplying butterfly matrices as follows: $\bm{B}_{d-1,1}(d)\bm{B}^\top_{d-1,2}(d)\cdots\bm{B}_{1,1}(d)\bm{B}^\top_{1,2}(d)$, where $\bm{B}_{i,j}(d),\forall i,\forall j$ are butterfly matrices.
\end{theorem}

Theorem~\ref{thm:exp_boft} inspires a direct generalization of BOFT. Specifically, the resultant orthogonal matrix can be formulated as $\thickmuskip=2mu \medmuskip=2mu \bm{R}^G(m_1,b_1,m_2,b_2,l)=\bm{R}_{l,1}(m_1,b_1)\bm{R}_{l,2}^T(m_2,b_2)\cdots\bm{R}_{1,1}(m_1,b_1)\bm{R}_{1,2}^T(m_2,b_2)$, with $\bm{R}_{i,j}^T(m,b)$ representing the orthogonal matrix element within BOFT. When setting $m_1=m_2=\log d$, $b_1=b_2=2$ and $l=d-1$, this construction $\bm{R}^G(m_1,b_1,m_2,b_2,l)$ can capture the entirety of the orthogonal group, a property also described as the kaleidoscope hierarchy~\cite{dao2019kaleidoscope}. Nevertheless, it is important to recognize that higher expressivity does not universally translate to improved downstream finetuning performance. Indeed, fully expressive models, such as those resulting from complete finetuning, can perform suboptimally in practice. Thus, the balance between expressiveness and regularization fundamentally governs the generalization potential in model finetuning. BOFT expands the set of finetunable parameters in a way that embeds structural assumptions, assisting practitioners in striking a more optimal equilibrium between expressivity and regularization.

\textbf{Spectral properties}. Compared to LoRA, orthogonal finetuning tends to achieve superior spectral characteristics, as it maintains the spectral norm of the original pretrained weight matrix $\bm{W}^0$ without alteration. This can be understood through the singular value decomposition: $\bm{W}^0=\bm{U}\bm{\Sigma}\bm{V}^\top$, where $\bm{U}$ and $\bm{V}$ are both orthogonal matrices, and $\bm{\Sigma}$ contains the singular values along its diagonal. When either OFT or BOFT is applied, they act by left-multiplying the weight matrix by another orthogonal matrix $\bm{R}$, resulting in adjusted weights $\bm{R}\bm{U}\bm{\Sigma}\bm{V}^\top$. This operation leaves the largest singular value—and therefore the spectral norm—of $\bm{W}^0$ intact. Preservation of this spectral property has been documented as enhancing both training stability and generalization ability~\cite{miyato2018spectral,yoshida2017spectral}. Additional relevant mathematical characteristics are provided in Appendix~\ref{app:sn_property}.

\textbf{Orthogonal finetuning as learning bilinear similarity}. The BOFT method can be reformulated as learning the bilinear similarity $\bm{w}^0_i\bm{R}\bm{x}$, where $\bm{w}^0_i$ denotes the $i$-th column (neuron) in $\bm{W}^0$. Under this perspective, BOFT is essentially learning a bilinear similarity matrix $\bm{R}$, subject to strict constraints—namely, that $\bm{R}$ must remain orthogonal. This formulation creates a direct link to the concepts of distance metric learning~\cite{xing2002distance} as well as the mathematical framework of bilinear forms~\cite{roman2005advanced}.

\textbf{Inductive bias and generalization in BOFT}. The matrix $\bm{R}(m,b)$ in BOFT is typically drawn from a structured subclass of the full orthogonal group, thereby narrowing the function class available for optimization and imposing a particular inductive bias. We posit that the unique inductive bias associated with butterfly factorization can be advantageous for generalization, given that it encodes patterns common to several foundational linear transforms~\cite{dao2019kaleidoscope}, including the discrete Fourier transform, discrete sine/cosine transform, and the Hadamard transform. Furthermore, the sparse matrix structure inherent in the BOFT parameterization might confer additional, implicit forms of inductive bias~\cite{gunasekar2017implicit,li2018algorithmic,liu2019neural}.

\textbf{Contrast with butterfly-based sparse training}. Numerous prior studies~\cite{dao2019learning,dao2019kaleidoscope,chen2022pixelated,dao2022monarch} have explored sparse training techniques grounded in the butterfly parameterization. Typically, these approaches directly express weight matrices in neural networks via butterfly parameterization and train models from the ground up. In \cite{dao2022monarch}, the authors take a different route: they finetune pretrained models by projecting weights onto a variant of butterfly matrices before optimizing the resultant components for specific downstream tasks. In contrast, BOFT introduces a substantially different finetuning procedure that applies layer-wise transformations to weights using shared matrix parameters.

\input{rebuttal-results}

\section{Concluding Remarks and Limitations}

We introduce Orthogonal Butterfly, a universal parameter-efficient finetuning technique for foundation models built on the butterfly matrix concept. The main innovation for improved parameter-efficiency is representing a dense orthogonal matrix as a product of several sparse orthogonal matrices. To streamline the search for appropriate matrix decompositions, we design a graphical information transmission framework. Within this framework, our results reveal that the butterfly structure serves as an effective mechanism for achieving our goal of sparse orthogonal matrix decompositions. We provide empirical results supporting the practicality of BOFT for finetuning across large language models, extensive vision systems, and text-to-image generators. Additionally, our findings highlight the robustness of BOFT as a broadly applicable finetuning method.

However, BOFT is not without drawbacks. Since it constructs the final orthogonal transformation as a sequence of orthogonal matrices, the associated training overhead is marginally higher than that of OFT. Addressing and minimizing BOFT's training time overhead constitutes ongoing work. Nevertheless, after completion of the finetuning phase, BOFT-generated orthogonal matrices can be collapsed into the pretrained model, introducing no extra latency at inference. Furthermore, it remains unclear if the butterfly architecture constitutes the most optimal method for information propagation. Our information transmission framework offers opportunities to take cues from other disciplines—particularly computer networking, where network topology efficiency is well-studied. We anticipate that future work could employ even more efficient network designs for constructing dense orthogonal matrices.